\section{From data summaries to random quantities}

In the previous chapter, a dataset was treated as a finite list of observed
values. If the observations were
\[
x_1,x_2,\ldots,x_n,
\]
then quantities such as the sample mean, sample variance, median, empirical
quantiles, or empirical cumulative distribution function were computed from that
particular list. Once the data had been observed, these quantities were ordinary
numbers or ordinary functions.

This chapter changes the point of view again. We now ask what happens before the
sample is observed. Before observation, the entries in the sample are not fixed
numbers. They are random variables. Consequently, any summary computed from them
is also random.

For example, after observing five values we might compute
\[
\bar{x}=2.8.
\]
This is a number. But before the five observations are known, the sample mean is
\[
\overline X=\frac{1}{5}(X_1+X_2+X_3+X_4+X_5),
\]
and this is a random variable. It depends on the random sample
\[
X_1,X_2,X_3,X_4,X_5.
\]
A different sample may produce a different value of \(\overline X\).

\begin{summarybox}
A statistic computed from observed data is a number. The same statistic computed
from a random sample, before the sample is observed, is a random variable.
\end{summarybox}

\subsection*{Two stages of the same object}

It is useful to distinguish two stages.

\begin{enumerate}
    \item Before sampling, the observations are random variables
    \[
    X_1,X_2,\ldots,X_n.
    \]
    A statistic such as
    \[
    T=g(X_1,X_2,\ldots,X_n)
    \]
    is also a random variable.

    \item After sampling, the observations become concrete values
    \[
    x_1,x_2,\ldots,x_n.
    \]
    The same rule gives the observed value
    \[
    t=g(x_1,x_2,\ldots,x_n).
    \]
\end{enumerate}

The formula may look almost identical in both cases, but the mathematical status
is different. Before observation, \(T\) has a probability law. After observation,
\(t\) is one realized value.

\begin{remarkbox}
The distinction between \(T\) and \(t\) is not a typographical detail. It is the
basic distinction between a random procedure and the numerical result produced by
one execution of that procedure.
\end{remarkbox}

\subsection*{A first example: the sample mean}

Suppose a population model is represented by a random variable \(X\) with
expectation \(\mu\). Imagine that we take a sample of size \(4\). Before the
sample is observed, we write
\[
X_1,X_2,X_3,X_4.
\]
The sample mean is then
\[
\overline X=\frac{X_1+X_2+X_3+X_4}{4}.
\]
This quantity is random because each term in the numerator is random.

If one particular sample gives
\[
2.1,\quad 3.4,\quad 1.8,\quad 4.7,
\]
then the observed sample mean is
\[
\bar{x}=\frac{2.1+3.4+1.8+4.7}{4}=3.0.
\]
Another sample of size \(4\) from the same model might give a different value,
for example \(\bar{x}=2.6\) or \(\bar{x}=3.3\). The randomness is not in the
arithmetic rule. The randomness is in the sample to which the rule is applied.

\begin{examplebox}
If the data have already been collected, the sample mean is a descriptive
summary of those data. If the data are still to be collected, the sample mean is
a random variable whose possible values depend on the sampling process.
\end{examplebox}

\subsection*{Why this change of viewpoint matters}

This change of viewpoint is the beginning of statistical inference. Descriptive
statistics tell us what happened in one observed dataset. Inference asks what can
be concluded about an underlying model, population, or parameter from that one
dataset.

To answer such questions responsibly, we must understand how much a statistic can
vary from sample to sample. For instance, if two samples have different means,
there are at least two possible explanations:
\begin{itemize}
    \item the underlying populations or mechanisms may genuinely be different;
    \item the observed difference may be a consequence of random sampling
    variation.
\end{itemize}

Sampling distributions provide the mathematical language for separating these
possibilities. They describe how a statistic behaves under repeated sampling from
a specified model.

\subsection*{Observed distribution versus sampling distribution}

A common source of confusion is the word ``distribution''. In this part of the
course, two different distributions must be kept separate.

The empirical distribution describes the observed data values
\[
x_1,x_2,\ldots,x_n.
\]
It tells us where the observations are located and how they are spread.

A sampling distribution describes the possible values of a statistic such as
\[
\overline X,\qquad \widehat p,\qquad S^2,\qquad \text{or}\qquad M,
\]
where \(\widehat p\) may denote a sample proportion, \(S^2\) a sample variance,
and \(M\) a sample median. It tells us how that statistic would vary if the
sampling procedure were repeated many times.

\begin{center}
\begin{tabular}{lll}
\toprule
Object & Level & Typical question \\
\midrule
Observed data \(x_1,\ldots,x_n\) & data level & What did we observe? \\
Empirical distribution & data level & How are the observed values distributed? \\
Statistic \(T=g(X_1,\ldots,X_n)\) & sampling level & How does the summary vary? \\
Sampling distribution of \(T\) & sampling level & What values of \(T\) are likely? \\
\bottomrule
\end{tabular}
\end{center}

The empirical distribution belongs to one dataset. The sampling distribution
belongs to a random procedure.

\subsection*{The guiding idea of the chapter}

The central idea of this chapter can now be stated in one sentence: every
statistic that is computed from a random sample has its own probability law.

This law may sometimes be found exactly. In other cases, it may be approximated
mathematically or explored by simulation. In all cases, the purpose is the same:
to understand the random variability of summaries computed from samples.

\begin{theorembox}
The main conceptual step in this chapter is to move from
\[
\text{data } \longrightarrow \text{ summaries of data}
\]
to
\[
\text{random samples } \longrightarrow \text{ random summaries}.
\]
This step is what makes statistical inference possible.
\end{theorembox}
